\chapter{\label{chapter:intro}Introduction}

Embedded systems are being widely adopted in all areas of human activity. Their increased performance during recent years has enabled a wide range of new applications. However, more processing power implies more functionalities and increased software complexity, which directly impacts the design constraints and goals of embedded systems \cite{sisu2014}. For example, the possibility for the final user to develop and download new applications, before exclusively found in general-purpose systems, can now be found in many modern embedded devices \cite{HeiserRole}. This scenario has forced designers and companies to adopt new strategies, such as increased use of software layers allowing for more flexible platforms capable of meeting timing, energy consumption, and time-to-market constraints. Although embedded systems are assuming many of the features of general-purpose systems, some differences remain. They present critical real-time constraints and frequently have resource constraints such as, limited battery, memory and processing power. 

The appearance of server farms in the 1990s and the increasing performance of the computer systems stimulated the adoption of virtualization. Virtualization allows for the possibility of decoupling the one-to-one correspondence between hardware and software. Thus, keeping several different operating systems on the same computer system, can drastically reduce the maintenance and energy costs \cite{muench2014}.  Different operating systems reside in a shared memory and are combined on a single or multi-core processor through a software stack called a virtual machine monitor or hypervisor. In recent years,  virtualization technology has quickly moved towards embedded systems motivated by the increasing processing power of the embedded processors and the increasing challenges to comply with their requirements. Although server virtualization is a well-known and mature technology widely applied for commercial use, embedded system virtualization is still being studied and developed. The main restriction to broader use of virtualization for embedded systems is that their requirements differ from server and enterprise systems \cite{Sandtrom2013}. Most notably, timing constraints and limited resources, like CPU and memory, are the main concerns around embedded systems virtualization. Thus, virtualization as deployed in the enterprise environment cannot be directly used in embedded systems. 

There are many studies proposing different techniques for embedded system virtualization. A common approach is to adapt hypervisors widely used in server virtualization to embedded systems. The main concern about general purpose hypervisors is the absence of suitable support for real-time and memory requirements incompatible with some embedded devices. Despite these limitations, researchers are continually working to improve the embedded support of open-source hypervisors for server virtualization, like XEN and KVM, as seen in \cite{Dall2014}, \cite{Dan2015}, \cite{seehwan2014} and \cite{avanzini2015}. On the other hand, the distinguished characteristics of the embedded systems have motivated the appearance of hypervisors specially designed for embedded virtualization. Among the goals for embedded hypervisors development, two of them are frequently addressed: to keep low memory requirements and some level of support for real-time applications. Additionally, the diversity of embedded systems and their applications encourages the development of many different embedded hypervisors, like SPUMONE \cite{spumone2008}, Xtratum \cite{trujillo2013}, M-Hypervisor \cite{rui2013}, Xvisor \cite{patel2015} and the L4 hypervisors family \cite{Xia2015}, \cite{Heiser2010}, \cite{pikeos}, \cite{yang2011}.

Intel and AMD manufactures provided hardware support for virtualization on the x86 architecture to address the continuous adoption of server virtualization. This was aimed  to simplify the virtualization software layer and to improve performance. Before hardware-assisted virtualization, hypervisors needed to implement a technique called para-virtualization that consists of modifying the operating system to be virtualized. These operating system modifications break the software compatibility and require additional engineering work along with licence and source-code access for proprietary software. With dedicated hardware for virtualization, the hypervisors could implement a technique called full-virtualization where no modifications are required on the target operating system. Thus, making it easier to support new operating systems, especially proprietary ones. Similar to the x86 architecture, the late adoption of hardware-assisted virtualization for embedded processors made para-virtualization rule the development of embedded virtualization. Nowadays, the main embedded processor manufacturers had already designed virtualization extensions for their processor families, like PowerPC \cite{powerpcvz}, ARM \cite{armvt} and MIPS \cite{mipsvz}. These gave hypervisor developers more design choices. Modern hardware-assisted virtualization now makes it possible to virtualize an operating system without any modification and keeps performance close to the level of non-virtualized systems. This is important to support legacy software. However, advanced hypervisor features like communication among virtual machines or shared devices may still require para-virtualization technique. 

This dissertation starts by addressing the different embedded virtualization approaches and discussing their drawbacks and advantages. It points to several features required for embedded virtualization that the current hypervisor approaches fail to support. Thus, it is proposed the hybrid virtualization model concept that combines full and para-virtualization in the same hypervisor. To support full-virtualization, the model considers the recent hardware-assisted virtualization adopted in some embedded processors. Para-virtualization is used for extended services, such as communication among virtual machines, and real-time services. A hypervisor implementation based on the proposed model was developed for the MIPS M5150 processor and evaluated in the SEAD-3 development board. The contribution of this Ph.D. research is to show how design choices based on hardware-assisted virtualization associated with embedded system characteristics result in a simpler hypervisor, while still improving overall performance of the whole system. 


\section{Hypothesis and Research Questions}

The aim of this work is to investigate the hypothesis that is possible to have an embedded lightweight hypervisor capable of supporting the major embedded virtualization requisites. Moreover, if considered that this hypothesis can be true, how important is the new hardware-assisted virtualization to accomplish it? Fundamental research questions associated with the hypothesis that guided this research are defined as follows:

%The aim of this Ph.D. research is to verify the hypothesis that the definition of an embedded virtualization model using hardware-assisted virtualization would result in a ligh\-tweight hypervisor capable of supporting the key required features for embedded virtualization. Additionally, it would present better responsiveness, footprint and performance than state-of-the-art hypervisors

\begin{enumerate}
  \item What is the state-of-the-art in embedded virtualization and how does it accomplish the needs of  embedded systems? This research question’s main objective is to study embedded virtualization in detail and identify where the current virtualization models fail to support virtualization for embedded systems.  Answering this research question will make it possible to understand the approaches that have already been tested, identify their limitations, and point out opportunities for improvements.

  \item Can a virtualization model embrace the major embedded systems’ needs and what features does it need to support to accomplish this goal? The objective of this research question is to determine the required embedded virtualization features and to propose a virtualization model capable of supporting these features. It is important to understand how different features, some of them contradictory, can be combined in the same model. The resulting virtualization model will guide the development of a new embedded hypervisor.
  
  \item Can an implementation from a virtualization model be trustworthy and how can hardware-assisted virtualization be used to achieve this goal? This research question addresses whether  a practical implementation can be constructed from the proposed virtualization model. Theoretical models can be hard to implement or an implementation may not achieve the expected results. 

\end{enumerate}


\section{Organization of This Dissertation}

The remainder of this dissertation is organized as follows:

\begin{itemize}
\item \textbf{Chapter \ref{chapter:background}} explains the basic concepts for the best understanding of the remainder of the text. It starts by defining embedded systems and their applications. Virtualization is defined along with a taxonomy of virtual machines. The different types of hypervisors are shown and the requirements for virtualizing a processor are explained. The technologies that make virtualization possible are also discussed. Moreover, two critical problems resulting from virtualization usage are presented.  Then, the relationship between virtualization and embedded systems is explained. Finally, it discusses the past, present and future possibilities for virtualization in embedded systems. 

\item \textbf{Chapter \ref{chapter:stateofart}} describes the state-of-the-art in embedded virtualization. The chapter starts by showing hypervisors designed for server virtualization. Then it presents several different authors’ proposals to adapt these hypervisors for embedded systems. Subsequently, different hypervisors developed specially for embedded systems are described.  A comparative study about the hypervisors is then presented. Finally, a Linux approach to improve real-time responsiveness is explained and compared to virtualization. Chapter \ref{chapter:stateofart} addresses research question (1).

\item \textbf{Chapter \ref{chapter:vtmodel}} presents the virtualization model proposed in this dissertation. First, it discusses the required features for a virtualization model for embedded systems. Then, the proposed embedded virtualization model is presented, and its different aspects are discussed. Finally, the model is compared to the Virtual Hellfire Hypervisor's model (VHH is a hypervisor previously developed in the Embedded Systems Group at PUC-RS). Chapter \ref{chapter:vtmodel} addresses research question (2).

\item \textbf{Chapter \ref{chapter:implementation}} describes the hypervisor implementation based on the virtualization model proposed. First, the processor supported by the hypervisor is briefly presented. The software architecture is discussed and the CPU virtualization strategy is explained.  Then the memory virtualization approach is exposed and the strategy for interrupt handling is described. The mechanism to allow communication among virtual machines is explained. Finally, the engineering effort to support a new guest and the implementation restrictions are discussed. Chapter \ref{chapter:implementation} partially addresses research question (3).

\item \textbf{Chapter \ref{chapter6:evaluation}} shows the hypervisor evaluation and practical results. The hypervisor memory requirements are determined and compared to other hypervisors. Next, the experience to port Linux and minor modifications to the Linux kernel to improve performance is described. The overall hypervisor overhead impact is then analyzed. The interrupt delivery strategy and the communication mechanism among virtual machines are also evaluated. Moreover, the real-time capabilities are measured. Lastly the overall results are discussed.  Chapter \ref{chapter6:evaluation} concludes research question (3).

\item \textbf{Chapter \ref{chapter:conclusions}} summarizes the dissertation and presents the concluding remarks. It restates the answers to the research questions, the main contributions, and presents possible directions for future work.
\end{itemize}